\documentclass[UTF8,12pt,a4paper,fontset=fandol]{ctexart}

% 支持从项目根目录或当前笔记目录编译。
\makeatletter
\def\input@path{{../}{../../}}
\makeatother
\usepackage{researchnotes}

\title{对角阵、单位阵与上三角阵}
\author{}
\date{}

\begin{document}
\maketitle

\section{主对角线与对角阵}

对角阵、单位阵和上三角阵，都是由元素所在位置来描述的特殊方阵。
以下设 $n$ 为正整数，$A=(a_{ij})\in\mathbb R^{n\times n}$ 为 $n$ 阶方阵，
$a_{ij}$ 表示第 $i$ 行、第 $j$ 列的元素。
从左上角到右下角的 $a_{11},a_{22},\ldots,a_{nn}$ 构成
\keyterm{主对角线}（main diagonal）。主对角线上满足 $i=j$，其上方满足 $i<j$，下方满足 $i>j$。

\begin{definition}[对角阵]
若主对角线以外的元素全部为 $0$，即 $i\ne j$ 时 $a_{ij}=0$，则称 $A$ 为
\keyterm{对角矩阵}（diagonal matrix），简称\keyterm{对角阵}。
\end{definition}

记号 $\operatorname{diag}(d_1,\ldots,d_n)$ 表示主对角线元素依次为 $d_1,\ldots,d_n$ 的对角阵。例如
\[
  D=\begin{pmatrix}2&0&0\\0&0&0\\0&0&-1\end{pmatrix}
   =\operatorname{diag}(2,0,-1).
\]
\keyterm{对角线上的元素可以为零}，也可以为正数或负数。
“只有主对角线上的元素可以非零”不等于“主对角线上的元素必须非零”。
因此，元素全为 $0$ 的方阵也是对角阵。

对角阵乘列向量，就是分别缩放各个坐标：若 $x=(x_1,\ldots,x_n)^{\mathsf T}$，则
$\operatorname{diag}(d_1,\ldots,d_n)x=(d_1x_1,\ldots,d_nx_n)^{\mathsf T}$，
其中 $\mathsf T$ 表示转置。例如，上述 $D$ 将 $(x_1,x_2,x_3)^{\mathsf T}$ 变为
$(2x_1,0,-x_3)^{\mathsf T}$，各坐标之间没有相加混合。

\section{单位阵：保持向量不变}

\begin{definition}[单位阵]
主对角线元素全部为 $1$、其余元素全部为 $0$ 的方阵，称为
\keyterm{单位矩阵}（identity matrix），简称\keyterm{单位阵}，记为 $I_n$。
因此 $I_n=\operatorname{diag}(1,\ldots,1)$，其中共有 $n$ 个 $1$。
\end{definition}

例如，二阶和三阶单位阵分别为
\[
  I_2=\begin{pmatrix}1&0\\0&1\end{pmatrix},\qquad
  I_3=\begin{pmatrix}1&0&0\\0&1&0\\0&0&1\end{pmatrix}.
\]
单位阵是特殊的对角阵，并不是所有元素都为 $1$ 的矩阵。
它把每个坐标乘以 $1$，所以对任意 $x\in\mathbb R^n$ 都有 $I_nx=x$。
更一般地，对 $B\in\mathbb R^{m\times n}$，有
\[
  I_mB=B=BI_n.
\]
左乘单位阵保留每一行，右乘单位阵保留每一列。
所以单位阵在矩阵乘法中起到实数 $1$ 在数的乘法中的作用；左右两侧单位阵的阶数由尺寸决定。

\section{上三角阵：主对角线下方全为零}

\begin{definition}[上三角阵]
若主对角线下方的元素全部为 $0$，即 $i>j$ 时 $a_{ij}=0$，则称 $A$ 为
\keyterm{上三角矩阵}（upper triangular matrix），简称\keyterm{上三角阵}。
\end{definition}

三阶上三角阵的一般形式及一个具体例子为
\[
  \begin{pmatrix}a_{11}&a_{12}&a_{13}\\0&a_{22}&a_{23}\\0&0&a_{33}\end{pmatrix},
  \qquad
  U=\begin{pmatrix}1&2&-1\\0&0&3\\0&0&4\end{pmatrix}.
\]
“上三角”表示非零元素只能出现在主对角线及其上方；这些位置也允许为零。
尤其是，主对角线上出现 $0$ 不影响它成为上三角阵，上述 $U$ 就是一个例子。

上三角结构使线性方程组便于从下向上求解。例如
\[
  \begin{pmatrix}1&2\\0&3\end{pmatrix}
  \begin{pmatrix}x_1\\x_2\end{pmatrix}
  =\begin{pmatrix}5\\6\end{pmatrix}
  \quad\Longleftrightarrow\quad
  \begin{cases}x_1+2x_2=5,\\3x_2=6.\end{cases}
\]
先由第二式得 $x_2=2$，再代入第一式得 $x_1=1$。
这种方法称为\keyterm{回代}（back substitution）。
一般的上三角方程组在主对角线元素都非零时，可依次回代得到唯一解；
若其中有零，则需要另行判断是否无解或有自由变量，不能直接除以该对角元素。

\section{三者的关系与判断方法}

\keyterm{单位阵一定是对角阵，对角阵一定是上三角阵}。
前者因为单位阵的非对角元素全为零，后者因为对角阵的主对角线下方必然全为零。
反过来一般不成立：$\begin{pmatrix}2&0\\0&2\end{pmatrix}$ 是对角阵，但不是单位阵；
$\begin{pmatrix}1&2\\0&1\end{pmatrix}$ 是上三角阵，但不是对角阵。
判断时，先确认矩阵是方阵，再检查必须为零的位置：下方全零即可判为上三角阵；
非对角位置全零即可判为对角阵；在此基础上，对角元素全为 $1$ 才是单位阵。

\end{document}
